\chapter{Polarisability of an ellipsoidal particle\label{App:pol_oblate_sphe}}

In this appendix we define the polarisability of an ellipsoidal particle. A detailed demonstration of the following expressions can be found in chapter (5) of Bohren and Huffman book \cite{bohren2008absorption}.\\
Let us consider a particle in vacuum, with semiaxes $a$, $b$ and $c$, like the one presented in Fig. (\ref{fig:app_ellip_part}).The particle is characterised by its dielectric tensor $\boldsymbol{\epsilon}$.\\

\begin{figure}[h!]
\begin{centering}
\includegraphics[width=7cm]{appendices/oblate.png} 
\par\end{centering}
\caption{Representation of a ellipsoidal particle. $a$, $b$ and $c$ are the particle's dimensions in the principal axes.\label{fig:app_ellip_part}}
\end{figure}
The static polarisability of an ellipsoidal particle is given by:

\begin{equation}
\boldsymbol{\alpha}_{0} = 4\pi abc \frac{\boldsymbol{\epsilon}-\mathbb{I}}{3\mathbb{I}+3\mathbb{L}\left(\boldsymbol{\epsilon}-\mathbb{I}\right)},
\end{equation}
being $\mathbb{L}$ the geometrical factors and defined by:

\begin{equation}
\mathbb{L} = \left[\begin{array}{ccc}
L_{1} & 0 & 0\\
0 & L_{2} & 0\\
0 & 0 & L_{3}
\end{array}\right].
\end{equation}

The geometrical factors can be calculated by considering the following expressions:
\begin{eqnarray}
L_{1} & = & \frac{abc}{2}\int^{\infty}_{0} \frac{dq}{\left(a^2+q\right)f\left(q\right)} \nonumber \\
L_{2} & = & \frac{abc}{2}\int^{\infty}_{0} \frac{dq}{\left(b^2+q\right)f\left(q\right)} \nonumber \\
1 & = & L_{1} + L_{2} + L_{3},
\end{eqnarray}
with $f\left(q\right) = \left\{\left(q+a^2\right)^2\left(q+b^2\right)^2\left(q+c^2\right)\right\}^{1/2}$.

An interesting case of ellipsoids are the spheroids, that have two axes of equal length, and consequently only one of the geotrical factors is independent.  Prolates and oblates ares generated by rotating an ellipse about its major and minor axis respectively. \\ In prolates we have $L_{2} = L_{3}$ and in oblates $L_{1} = L_{2}$.
Thus for a prolate spheroid ($b = c$):

\begin{eqnarray}
L_{1} & = & \frac{1-e^2}{e^2}\left(-1+\frac{1}{2e}\ln\frac{1+e}{1-e}\right)\\
 e^2& = &1-\frac{b^2}{a^2},
\end{eqnarray}

and for the oblate spheroid ($a = b$):

\begin{eqnarray}
L_{1} &=& \frac{g\left(e\right)}{2e^2}\left[\frac{\pi}{2} - \tan^{-1}g\left(e\right)\right]-\frac{g^2\left(e\right)}{2}\\
g\left(e\right) & = &\left(\frac{1-e^2}{e^2}\right)^{1/2}\\ 
e^2& = &1-\frac{c^2}{a^2}.
\end{eqnarray}

Both cases, when $e = 0$, result in a sphere.